\begin{abstract}
Equatorial Plasma Bubbles (EPBs) produce ionospheric irregularities that cause radio scintillation and degrade GNSS and satellite communication links at low latitudes. Optical monitoring of EPBs still depends largely on visual inspection of $630.0\text{ nm}$ all-sky airglow images, a slow task made harder by moonlight, horizon brightening, and clouds. We describe a two-stage pipeline for automatic EPB detection. The first stage checks sky transparency by counting star pixels and measuring the mean sky brightness inside the field of view, and discards cloud-covered frames. The second stage removes the diffuse background, enhances local contrast with CLAHE, extracts features with Two-Dimensional Principal Component Analysis (2DPCA), and classifies them with a stacking ensemble of SVM, XGBoost, and Random Forest. On $572$ labeled images from three EMBRACE observatories in Brazil, the classifier reached $97.20\%$ accuracy and $0.998$ ROC-AUC under stratified 5-fold cross-validation, and $95.08\%$ accuracy when whole $30$-minute windows were held out. Without background normalization, recall on an observatory left out of training fell from $100\%$ to $0\%$, which indicates that the model then learns station brightness rather than plasma structure. On an overcast night excluded from training, the first stage flagged $79\%$ of the frames as cloudy and cut false alarms by $78.7$--$90.2\%$. SHAP maps projected back to pixel space show that the predictions follow the north--south depletion bands.

\keywords{Equatorial Plasma Bubbles \and Machine Learning \and Space Weather \and All-Sky Imager \and 2DPCA \and Stacking Ensemble \and Explainable AI.}
\end{abstract}
